% STARTER DRAFT for Emmett -- written with Claude Code on 2026-10-20 at Patrick's
% request. Not included in main.tex. Use, rewrite or ignore; see drafts/README.md.

% ---------------------------------------------------------------------
% -> sections/related_work.tex, interpretability paragraph (shared with Conor)
% Sources: journal/2026-10-07-circuit-tracing.md; the two Transformer Circuits pages
% (section list checked on 2026-10-20).
% ---------------------------------------------------------------------
\textbf{Looking inside the model.} Circuit tracing replaces a model's MLPs with
transcoders and builds attribution graphs that show which features drive a
prediction~\cite{ameisen2025}. Lindsey et al.\ apply it to case studies that include
addition and chain-of-thought faithfulness, where they contrast a faithful
computation with cases where the reasoning is made up or worked backwards from a
suggested answer~\cite{lindsey2025}.
% CHECK (Emmett): you were going to read these sections ("Still open" in your
% 2026-10-07 entry) -- confirm the description before citing.

% ---------------------------------------------------------------------
% -> sections/method.tex, subsection "Planted working and copy-vs-use"
% Sources: journal/2026-10-07-circuit-tracing.md; journal/2026-10-15.md (decision 4)
% ---------------------------------------------------------------------
Circuit-tracing transcoders for Qwen3-1.7B did not fit in free Colab memory, so these
tests use Qwen3-0.6B with short prompts of the form ``Question: $a+b=?$ / Working:
\dots / Final answer:''. We record the probability of the first answer digit (digits
are single tokens in Qwen), comparing a wrong working, a correct working and no
working, on 8 hand-made addition prompts.

A wrong working that the model follows could be used or simply copied. The copy-vs-use
control separates the two with prompts whose answer needs one more step after the
working. For example, $(12+7)\times 2$ with working $12+7=21$: using the working gives
42, copying it gives 21, and the true answer is 38, so the first digits differ.
Prompts: \textbf{TBD}.
% TODO (Emmett): n, how prompts were built, any change from the 2026-10-15 design.

% ---------------------------------------------------------------------
% -> sections/results.tex, subsection "Planted working and copy-vs-use"
% Sources: outputs/circuits/hand_made_behavioural.csv, hand_made_baseline.csv
% (Table II already summarises them -- check it); journal/2026-10-07-circuit-tracing.md
% ---------------------------------------------------------------------
With no working, Qwen3-0.6B gives the true answer's first digit with probability
0.68--0.83 (mean 0.75), and with a correct working 0.99 on all 8 prompts
(Table~\ref{tab:planted}). With a wrong working, the first digit follows the working
with probability 0.94--0.99, and the true digit falls to at most 0.013. A wrong working
therefore overrides the model's own default answer. These are 8 hand-picked prompts
on a small model, so they demonstrate the effect rather than measure a rate.

Copy-vs-use: \textbf{TBD} (Table~\ref{tab:copyuse}).
% Optional figure: outputs/circuits/hand_12plus7_working21.png, as an illustration
% only -- the graph's edges have not been inspected, so it is not evidence of dependence.
